%% Fasit til kortsett fra mattelab.no – laget automatisk av bygg.py.
%% Kompiler med lualatex.
\documentclass[11pt,a4paper]{article}
\usepackage[norsk]{babel}
\usepackage[a4paper,margin=1.2cm,bottom=1.7cm]{geometry}
\usepackage{fontspec}
\usepackage{unicode-math}
\setmainfont{TeX Gyre Pagella}
\setsansfont{TeX Gyre Heros}[Scale=MatchLowercase]
\setmathfont{TeX Gyre Pagella Math}
\usepackage{xcolor}
\definecolor{blekk}{HTML}{16213A}
\definecolor{dempet}{HTML}{566077}
\definecolor{m1}{HTML}{2F4AB0}
\usepackage{tikz}
\usepackage{enumitem}
\usepackage{qrcode}
\usepackage{fancyhdr}
\usepackage[hidelinks]{hyperref}
\pdfvariable suppressoptionalinfo \numexpr 1+2+4+8+16+32+64+128+256+512\relax
\setlength{\parindent}{0pt}
\setlength{\parskip}{0pt}
\pagestyle{fancy}
\fancyhf{}
\renewcommand{\headrulewidth}{0pt}
\fancyfoot[L]{\sffamily\footnotesize\color{dempet}mattelab.no · Kontinuerlig, knekk eller sprang? · kortsett K13 · fasit}
\fancyfoot[R]{\sffamily\footnotesize\color{dempet}Side \thepage}
\begin{document}
\begin{minipage}[t]{0.8\linewidth}\vspace{0pt}
{\sffamily\small\color{dempet}Fasit for læreren · kortsett K13 · R1}\par\vspace{0.25em}
{\fontsize{20}{24}\selectfont\bfseries Kontinuerlig, knekk eller sprang?}\par\vspace{0.5em}
Sorter kortene etter hva som skjer med funksjonen i x = 1. Er den kontinuerlig og deriverbar der, kontinuerlig med en knekk, eller ikke kontinuerlig? Begrunn med grenseverdier eller med grafen.\par\vspace{0.4em}{\small\color{dempet}Kortene er nummerert i den rekkefølgen de står i utskriften.}
\end{minipage}\hfill
\begin{minipage}[t]{0.15\linewidth}\vspace{0pt}\raggedleft
\qrcode[height=1.9cm]{https://mattelab.no/varierende/kontinuitet-og-deriverbarhet.html}
\end{minipage}\par\vspace{0.8em}

\subsection*{Kontinuerlig og deriverbar i x = 1}
\begin{itemize}[leftmargin=1.2em, itemsep=0.35em]
\item \textsf{\small K13-05}\quad \emph{Figur:} Parabel med bunnpunkt i (1, 1).\\{\small\color{dempet}Glatt graf med vannrett tangent i $x = 1$. Da er $f'(1) = 0$.}
\item \textsf{\small K13-06}\quad $f(x) = \begin{cases} x^2 & x < 1 \\ 2x - 1 & x \geq 1 \end{cases}$\\{\small\color{dempet}Begge delene gir $1$ i $x = 1$, og begge har stigningstall $2$ der. Grafen er glatt, selv om forskriften er delt.}
\item \textsf{\small K13-09}\quad $f(x) = \begin{cases} 2x & x \leq 1 \\ x^2 + 1 & x > 1 \end{cases}$\\{\small\color{dempet}Begge delene gir $2$ i $x = 1$, og begge har stigningstall $2$ der.}
\item \textsf{\small K13-10}\quad $f(x) = x^3 - 3x$\\{\small\color{dempet}Polynomfunksjoner er kontinuerlige og deriverbare overalt.}
\item \textsf{\small K13-13}\quad \emph{Figur:} Graf som ser ut som én glatt kurve gjennom (1, 1). Til venstre en parabel, til høyre en rett linje.\\{\small\color{dempet}Dette er $x^2$ til venstre og $2x - 1$ til høyre. Linja er tangenten til parabelen, så grafen er glatt.}
\end{itemize}
\subsection*{Kontinuerlig, men ikke deriverbar i x = 1}
\begin{itemize}[leftmargin=1.2em, itemsep=0.35em]
\item \textsf{\small K13-02}\quad \emph{Figur:} Graf formet som en V, med spissen i (1, 3).\\{\small\color{dempet}Spiss i $(1, 3)$: Grafen henger sammen, men har ingen tangent der.}
\item \textsf{\small K13-03}\quad $f(x) = |x - 1|$\\{\small\color{dempet}Grafen er en V med spissen i $(1, 0)$. Stigningstallet er $-1$ fra venstre og $1$ fra høyre.}
\item \textsf{\small K13-04}\quad $f(x) = \begin{cases} x^2 & x < 1 \\ 3x - 2 & x \geq 1 \end{cases}$\\{\small\color{dempet}Begge delene gir $1$ i $x = 1$, så f er kontinuerlig. Men stigningstallet er $2$ fra venstre og $3$ fra høyre.}
\item \textsf{\small K13-08}\quad Du betaler 30 \% skatt av det du tjener over 1 million. f(x) er skatten når du tjener x millioner.\\{\small\color{dempet}Skatten er $0$ opp til $1$ million og øker jevnt etter det. Grafen henger sammen, men har en knekk i $x = 1$.}
\end{itemize}
\subsection*{Ikke kontinuerlig i x = 1}
\begin{itemize}[leftmargin=1.2em, itemsep=0.35em]
\item \textsf{\small K13-01}\quad En parkering koster 30 kr per påbegynte time. f(t) er prisen etter t timer.\\{\small\color{dempet}Prisen hopper fra $30$ kr til $60$ kr når du passerer én time. Slike trappefunksjoner er vanlige i priser og avgifter.}
\item \textsf{\small K13-07}\quad $f(x) = \begin{cases} x & x \neq 1 \\ 2 & x = 1 \end{cases}$\\{\small\color{dempet}Grenseverdien er $1$, men $f(1) = 2$. Ett punkt er flyttet, og det er nok til at f ikke er kontinuerlig.}
\item \textsf{\small K13-11}\quad \emph{Figur:} Graf som går opp til et fylt punkt i (1, 2) og fortsetter fra en åpen sirkel i (1, 4).\\{\small\color{dempet}Grafen gjør et sprang fra $2$ til $4$. Den fylte prikken viser at $f(1) = 2$.}
\item \textsf{\small K13-12}\quad $f(x) = \begin{cases} x^2 & x < 1 \\ x + 1 & x \geq 1 \end{cases}$\\{\small\color{dempet}Fra venstre går f mot $1$, men $f(1) = 2$. Grafen gjør et sprang.}
\end{itemize}
\end{document}
